\documentclass[11pt,a4paper]{article}
\usepackage[T1]{fontenc}
\usepackage{lmodern,amsmath,amssymb,amsthm,mathtools,booktabs,microtype}
\usepackage[margin=27mm]{geometry}
\usepackage{xcolor}
\definecolor{ink}{HTML}{164567}
\usepackage[colorlinks=true,linkcolor=ink,citecolor=ink,urlcolor=ink]{hyperref}
\hypersetup{pdftitle={The sixth-order cost of spectral asymmetry in shortest positive filters},pdfauthor={Research note prepared for Vladimir Reshetnikov}}
\newtheorem{theorem}{Theorem}[section]
\newtheorem{lemma}[theorem]{Lemma}
\newtheorem{corollary}[theorem]{Corollary}
\newtheorem{proposition}[theorem]{Proposition}
\theoremstyle{remark}\newtheorem{remark}[theorem]{Remark}
\newcommand{\R}{\mathbb R}\newcommand{\C}{\mathbb C}
\newcommand{\dd}{\,\mathrm d}\newcommand{\ii}{\mathrm i}
\newcommand{\supp}{\operatorname{supp}}\newcommand{\diam}{\operatorname{diam}}
\newcommand{\ord}{\operatorname{ord}}\newcommand{\Res}{\operatorname{Res}}
\newcommand{\coeff}[2]{[#1^{#2}]}
\setlength{\emergencystretch}{1.5em}
\allowdisplaybreaks[2]
\title{The sixth-order cost of spectral asymmetry\\in shortest positive filters}
\author{Research note prepared for Vladimir Reshetnikov}
\date{27 November 2026}
\begin{document}
\maketitle
\begin{abstract}
How short can a positive average be if it must cancel a cluster of oscillatory modes exactly? For a fixed number $N\ge3$ of nearby positive frequencies, with repeated entries prescribing repeated zeros, let $\kappa$ be their mean and $\delta_j=(\lambda_j-\kappa)/\kappa$. We prove that the minimum support diameter is
\[
 \frac{N\pi}{\kappa}
 +\frac{N\pi^5}{180(N-1)(N-2)\kappa}
       \left(\sum_j\delta_j^3\right)^2
 +O_N\!\left(\frac{\max_j|\delta_j|^8}{\kappa}\right).
\]
Thus the first five orders of spectral detuning do not change the optimum after the mean is fixed. The first possible change is sixth order, with an explicit positive coefficient determined only by the squared third centered moment. For three or four target occurrences, the mean-based width is an exact local lower bound, with equality precisely for a reflection-symmetric multiset. The result answers the first-next-term question in the supplied ProveIt article and sharpens its cubic remainder to the optimal sixth order. We give a self-contained proof, including analytic dependence through repeated roots and an exact second-variation calculation. The filter theorem is unconditional; its application to Fourier expansions of prime errors uses the repository's Dirichlet GRH input.
\end{abstract}

\section{The theorem and its significance}
For a compactly supported Borel probability measure $\mu$ on $\R$ write
\[
 \Phi_\mu(\omega)=\int_{\R}e^{\ii\omega u}\dd\mu(u).
\]
This is an entire function. A multiset $\boldsymbol\lambda=(\lambda_1,\ldots,\lambda_N)$ of positive numbers prescribes the constraints
\begin{equation}\label{eq:constraints}
 \ord_{\lambda}\Phi_\mu\ge\#\{j:\lambda_j=\lambda\}.
\end{equation}
Define $B(\boldsymbol\lambda)$ to be the infimum of $\diam\supp\mu$ over measures satisfying \eqref{eq:constraints}. Translating $\mu$ does not affect the constraints or its diameter. Rescaling the frequencies gives
\begin{equation}\label{eq:scale}
 B(c\boldsymbol\lambda)=c^{-1}B(\boldsymbol\lambda)\qquad(c>0).
\end{equation}
The feasible set is nonempty: convolve the two-point filters with atoms at $0$ and $\pi/\lambda_j$, each of mass $1/2$. Our optimum permits every probability measure, including continuous measures; the number of atoms is a conclusion.

\begin{theorem}[The sixth-order law]\label{thm:main}
Fix $N\ge3$. There are $\epsilon_N>0$ and $C_N<\infty$ with the following property. Put
\[
 \kappa=\frac1N\sum_{j=1}^N\lambda_j,\qquad
 \delta_j=\frac{\lambda_j-\kappa}{\kappa},\qquad
 d=\max_j|\delta_j|,\qquad p_3=\sum_j\delta_j^3.
\]
If $d<\epsilon_N$, the minimum is attained by a unique measure up to translation. With left endpoint $0$, it has exactly $N+1$ positive atoms. Moreover
\begin{equation}\label{eq:main}
 \left|B(\boldsymbol\lambda)-\frac{N\pi}{\kappa}
 -\frac{N\pi^5}{180(N-1)(N-2)\kappa}p_3^2\right|
 \le\frac{C_Nd^8}{\kappa}.
\end{equation}
The exponent six is optimal: along every fixed centered direction with $p_3\ne0$, the excess over $N\pi/\kappa$ has a strictly positive sixth-order leading term.
\end{theorem}

\begin{theorem}[Exact local lower bound for three or four occurrences]\label{thm:equality}
For $N=3$ or $N=4$, and sufficiently small $d$ as above,
\[
 B(\boldsymbol\lambda)\ge\frac{N\pi}{\kappa}.
\]
Equality holds if and only if the multiset $\{\lambda_j-\kappa\}$ is invariant under negation, counting multiplicity. For three occurrences this means that one frequency equals the mean and the other two are equidistant from it. For four occurrences it means that the frequencies form two reflected pairs.
\end{theorem}

Here ``third centered moment'' means the unnormalized sum $p_3$, not a standardized statistical skewness. In unnormalized frequency units, the correction is
\[
 \frac{N\pi^5}{180(N-1)(N-2)\kappa^7}
      \left(\sum_j(\lambda_j-\kappa)^3\right)^2.
\]
For $N\ge5$, \eqref{eq:main} does not by itself prove a local lower bound in directions with $p_3=0$: higher odd moments can then govern the first nonzero term.

\paragraph{What is new relative to the source.}
We use the supplied repository at commit
\texttt{00cadb961bbe77c74e5d082c4948f8884d58a970}. In its positive-filter section \cite{ProveIt}, Theorem \texttt{fil:thm:split} proves a unique analytic optimum near a repeated frequency and gives
$B=N\pi/\kappa+O_N(d^3/\kappa)$. Theorem \texttt{fil:thm:differential} gives the first variation at a reflected spectrum. The open questions explicitly ask for the first nonzero term beyond $N\pi/\kappa$ and its sign for nonsymmetric clusters. Our contributions are the squared-odd-coordinate structure, the explicit sixth-order term, the eighth-order remainder, and the exact three- and four-occurrence equality criterion. Sections~\ref{sec:analytic}--\ref{sec:mirror} re-prove the local inputs needed here, so the mathematical proof does not depend on running the vendored Lean project or accepting the surrounding research claims.

Positive input shaping is an established subject; two-mode constructions and their dependence on mode ratios appear in \cite{Crain1996,Singhose1997}. We do not claim to introduce positive filters or repeated-frequency cancellation. A focused search did not locate the coefficient in \eqref{eq:main}; that is evidence of novelty relative to the checked sources, not a proof of global bibliographic priority.

\section{An analytic optimum in polynomial coordinates}\label{sec:analytic}
We normalize $\kappa=1$. Let $e_k$ be the elementary symmetric functions of $\delta_1,\ldots,\delta_N$, and set
\begin{equation}\label{eq:poly}
 Q_{\boldsymbol e}(y)=y^N-e_1y^{N-1}+e_2y^{N-2}-\cdots+(-1)^Ne_N.
\end{equation}
The mean constraint is $e_1=0$. It will be useful to allow the real coefficients $e_k$ to vary independently, even when the roots of $Q_{\boldsymbol e}$ are not all real.

\begin{lemma}[Analytic interpolation remainder]\label{lem:remainder}
Fix a small circle $C$ about $0$, and take $\boldsymbol e$ close enough to $0$ that all roots of $Q_{\boldsymbol e}$ lie inside $C$. For an entire function $f$, define
\begin{equation}\label{eq:remainder}
 \mathcal R_{\boldsymbol e}f(y)=\frac1{2\pi\ii}
 \oint_C\frac{f(1+z)}{Q_{\boldsymbol e}(z)}
       \frac{Q_{\boldsymbol e}(z)-Q_{\boldsymbol e}(y)}{z-y}\dd z.
\end{equation}
It is a polynomial in $y$ of degree below $N$, analytic in the coefficients $\boldsymbol e$. It has the same values and derivatives as $f(1+y)$ at the roots of $Q_{\boldsymbol e}$ through their prescribed multiplicities. At $\boldsymbol e=0$ it is
$\sum_{k=0}^{N-1}f^{(k)}(1)y^k/k!$.
\end{lemma}
\begin{proof}
The quotient in \eqref{eq:remainder} is polynomial of degree $N-1$ in $y$. Cauchy's formula rewrites the integral as
\[
 f(1+y)-Q_{\boldsymbol e}(y)\frac1{2\pi\ii}
       \oint_C\frac{f(1+z)}{Q_{\boldsymbol e}(z)(z-y)}\dd z
\]
for $y$ inside $C$. This proves the interpolation assertion, including repeated roots. The denominator on the fixed contour is nonzero in a coefficient neighborhood, which proves analyticity. The last assertion is Taylor's formula.
\end{proof}

\begin{proposition}[The local analytic branch and global optimality]\label{prop:branch}
For all real $\boldsymbol e$ near $0$, there are real analytic positive weights $w_j(\boldsymbol e)$ and ordered nodes
\[
 0=u_0(\boldsymbol e)<u_1(\boldsymbol e)<\cdots<u_N(\boldsymbol e)
\]
satisfying $\sum_jw_j=1$ and
$\mathcal R_{\boldsymbol e}\Phi_{\boldsymbol e}=0$, where
$\Phi_{\boldsymbol e}(\omega)=\sum_jw_j e^{\ii\omega u_j}$.
At $\boldsymbol e=0$,
\begin{equation}\label{eq:base}
 u_j=j\pi,\qquad w_j=2^{-N}\binom Nj,\qquad
 \Phi_0(\omega)=2^{-N}(1+e^{\ii\pi\omega})^N.
\end{equation}
When the roots $\delta_j$ of $Q_{\boldsymbol e}$ are real, $L(\boldsymbol e)=u_N(\boldsymbol e)$ is the global minimum width for the frequencies $1+\delta_j$, and the minimizing measure is unique after fixing its left endpoint.
\end{proposition}
\begin{proof}
There are $2N+1$ real equations: the real and imaginary parts of the $N$ remainder coefficients, and total mass one. There are $2N+1$ real unknowns, the $N+1$ weights and the $N$ nodes other than $u_0=0$. At the base these equations are equivalent to
\[
 \sum_jw_j u_j^k e^{\ii u_j}=0\quad(0\le k<N),\qquad \sum_jw_j=1.
\]
In a vector in the kernel of their Jacobian, the imaginary parts are
\[
 \sum_{j=1}^N(-1)^jw_j(j\pi)^k\,\dot u_j=0\quad(0\le k<N).
\]
The Vandermonde matrix is invertible, so every $\dot u_j=0$. The remaining real moment equations have a one-dimensional weight kernel, spanned by the binomial weights; the mass equation removes it. Thus the full Jacobian is invertible. The real analytic implicit function theorem, and strict positivity and order at the base, give the asserted branch in independent coefficient coordinates.

For optimality, let $Y_{\boldsymbol e}(u)\in\R^{2N}$ consist of the real and imaginary remainder coefficients of $\omega\mapsto e^{\ii\omega u}$. The vectors
\[
 Y(u_j)\ (0\le j\le N),\qquad Y'(u_j)\ (1\le j<N)
\]
span a hyperplane. Indeed the weights give one relation among the first $N+1$ vectors; the invertibility of the full mass-and-moment Jacobian implies that there is no second relation and that adding $Y'(u_N)$ gives rank $2N$. Let its analytic normal $v_{\boldsymbol e}$ be normalized by
$v_{\boldsymbol e}\cdot Y'(u_N)=-1$, and set $g_{\boldsymbol e}(u)=v_{\boldsymbol e}\cdot Y_{\boldsymbol e}(u)$.
It vanishes at all nodes and has zero derivative at every interior node. At the base it is a positive multiple of
\[
 (-1)^{N-1}\sin u\prod_{j=1}^{N-1}(u-j\pi).
\]
This function is strictly positive between nodes, has positive second derivatives at interior nodes, and has simple zeros of the appropriate sign at the endpoints. Those strict properties persist in a smaller neighborhood. Compactness away from small node neighborhoods gives
$g_{\boldsymbol e}\ge0$ on $[0,L(\boldsymbol e)]$, with precisely the specified nodes as zeros.

Every feasible measure supported in that interval integrates $g_{\boldsymbol e}$ to zero, so is supported on its nodes. The mass-and-moment columns are independent, again by the full Jacobian, and therefore fix its weights. The positive endpoint masses force its diameter to equal $L$. A competitor of smaller diameter can be translated into the interval; translation preserves zero orders. Hence no smaller competitor exists, and equality is unique.
\end{proof}

\begin{remark}[Why complex auxiliary roots are legitimate]
Only real roots are used to state the filtering theorem. Independent coefficient coordinates in Proposition~\ref{prop:branch} are an analytic continuation of its moment equations. This allows us to compute a Taylor coefficient along $e_3=t$, other $e_k=0$, although $y^N-ty^{N-3}$ generally has nonreal roots. No claim that this auxiliary path is a real frequency cluster is needed.
\end{remark}

\section{Mirror stationarity and the missing five orders}\label{sec:mirror}
Here a centered multiset is \emph{reflected} if it is invariant under $\delta\mapsto-\delta$.

\begin{lemma}[Exact width and differential at a reflected spectrum]\label{lem:mirror}
For sufficiently small real reflected $\delta_j$, $L=N\pi$. At every such configuration, a variation of its real frequencies satisfies
\begin{equation}\label{eq:differential}
 \dot L=-\pi\sum_j\dot\lambda_j.
\end{equation}
\end{lemma}
\begin{proof}
Put $r_j=e^{\ii\pi\lambda_j}$ and $P(z)=\prod_j(z-r_j)$. Reflected pairs give factors
$z^2+2\cos(\pi\delta_j)z+1$, and central entries give $z+1$. All coefficients are strictly positive for sufficiently small $\delta_j$. The coefficients of $P/P(1)$ on the nodes $j\pi$ satisfy the required zeros, so uniqueness in Proposition~\ref{prop:branch} gives $L=N\pi$.

First suppose the roots are distinct. Write the first variation of the transform as
\[
 \Phi_1(\omega)=A(z)+\ii\omega B(z),\qquad z=e^{\ii\pi\omega},
\]
where $A,B$ are real polynomials of degree at most $N$, $A(1)=0$, and $B(0)=0$. If $j^*$ is the partner with $\lambda_{j^*}=2-\lambda_j$, subtracting the conjugate of the varied root equation for $j^*$ from the equation for $j$ gives
\[
 B(r_j)=-\frac{\pi}{2P(1)}r_jP'(r_j)
               (\dot\lambda_j+\dot\lambda_{j^*}).
\]
Lagrange interpolation of $B(z)/z$ therefore gives
\[
 \frac{B(z)}z=-\frac{\pi}{2P(1)}
  \sum_j(\dot\lambda_j+\dot\lambda_{j^*})\frac{P(z)}{z-r_j}.
\]
Its leading coefficient is $-\pi\sum_j\dot\lambda_j/P(1)$. It is also $w_N\dot L$, where $w_N=1/P(1)$, proving \eqref{eq:differential}. Continuity gives the formula at repeated configurations.
\end{proof}

\begin{proposition}[A square of the odd-coordinate ideal]\label{prop:ideal}
On the coefficient hyperplane $e_1=0$, there are analytic functions $G_{ab}$ near $0$ such that
\begin{equation}\label{eq:ideal}
 L(\boldsymbol e)-N\pi=
  \sum_{\substack{3\le a,b\le N\\a,b\ {
m odd}}}
     e_ae_bG_{ab}(\boldsymbol e).
\end{equation}
Consequently, for a constant $A_N$,
\begin{equation}\label{eq:structure}
 L(\boldsymbol e)=N\pi+A_Ne_3^2+O_N(d^8).
\end{equation}
\end{proposition}
\begin{proof}
The odd coordinates $e_3,e_5,\ldots$ vanish precisely when the centered root polynomial is even (if $N$ is even) or $y$ times an even polynomial (if $N$ is odd). Reflected real multisets with distinct entries yield a nonempty open subset of the even-coordinate subspace: use distinct positive squared roots as parameters and their nonzero Vandermonde determinant.

On this open subset Lemma~\ref{lem:mirror} gives $L=N\pi$. It also gives every derivative of $L$ in the coordinates $e_2,\ldots,e_N$ equal to zero, since the root-to-coefficient Jacobian at distinct roots is invertible and these variations keep $e_1$ fixed. The real analytic identity theorem extends the value identity and the vanishing odd derivatives to a neighborhood in the even-coordinate subspace. Taylor's formula in the odd coordinates now gives \eqref{eq:ideal}.

For actual roots with $|\delta_j|\le d$, $e_k=O_N(d^k)$. The lowest possible weighted degree in \eqref{eq:ideal} is six, from $e_3^2$. Its coefficient can vary first by a coordinate of weight two, since $e_1=0$. Every other product has weight at least eight. This proves \eqref{eq:structure}; in particular no term of weighted degree seven occurs.
\end{proof}

\section{Exact computation of the sixth-order coefficient}\label{sec:coefficient}
We compute the coefficient in \eqref{eq:structure} by varying the independent coordinate $e_3=t$ and setting all other coordinates to zero:
\begin{equation}\label{eq:probe}
 Q_t(y)=y^N-ty^{N-3},\qquad L(t)=N\pi+A_Nt^2+O(t^3).
\end{equation}
Proposition~\ref{prop:branch} justifies ordinary convergent Taylor expansions along this coefficient path. Write
\[
 w_j(t)=w_j+ta_j+t^2a_j^{(2)}+O(t^3),\qquad
 u_j(t)=j\pi+tv_j+t^2v_j^{(2)}+O(t^3),
\]
with $w_j=2^{-N}\binom Nj$. Anchoring gives $v_0=v_0^{(2)}=0$.

\subsection{First variation}
Set $z=e^{\ii\pi\omega}$ and write
\[
 \Phi_1(\omega)=A(z)+\ii\omega B(z),\qquad
 A(z)=\sum_ja_jz^j,\quad B(z)=\sum_jw_jv_jz^j.
\]
In the next calculations put $x=z+1$, $\ell=\log(1-x)$, and $y=\ell/(\ii\pi)$; thus at $\omega=1+y$ or $-1+y$ the same $z=-e^{\ii\pi y}$ occurs. Congruences below are convergent power-series congruences modulo $y^N$, equivalently $x^N$.

Divisibility of $\Phi_t(1+y)$ by $Q_t(y)$ gives
\[
 \Phi_1(1+y)\equiv-\frac{\Phi_0}{y^3}.
\]
Reality of the weights and nodes gives divisibility at the negative cluster by $y^N+ty^{N-3}$, and hence
$\Phi_1(-1+y)\equiv\Phi_0/y^3$.
Subtracting and averaging the equations gives
\begin{equation}\label{eq:ABcong}
 B\equiv\frac{\ii\Phi_0}{y^3},\qquad
 A\equiv\frac{\Phi_0}{y^2},\qquad \Phi_0=2^{-N}x^N.
\end{equation}
Mass conservation is $A(1)=0$, and anchoring is $B(0)=0$. Expanding $\log(1-x)$ in \eqref{eq:ABcong}, and using these conditions to fix the possible multiples of $x^N$, yields
\begin{align}
 A(z)&=-\frac{\pi^2}{2^{N+2}}(z-1)^2(z+1)^{N-2},\label{eq:A}\\
 B(z)&=\frac{\pi^3}{2^{N+1}}z(1-z)(z+1)^{N-3}.\label{eq:B}
\end{align}
For example, the relevant reciprocal-log expansions are
$\ell^{-2}=x^{-2}(1-x+O(x^2))$ and
$\ell^{-3}=-x^{-3}(1-3x/2+x^2/2+O(x^3))$.
Equation~\eqref{eq:B} has degree below $N$, so $v_N=0$. Its coefficient at $z^j$, divided by $w_j$, is
\begin{equation}\label{eq:v}
 v_j=\frac{\pi^3j(N-j)(N-2j)}{2N(N-1)(N-2)}\qquad(0\le j\le N).
\end{equation}

\subsection{Second variation}
At order $t^2$ put
\[
 \Phi_2(\omega)=A_2(z)+\ii\omega B_2(z)-\frac{\omega^2}2 C(z),
\]
where
\[
 B_2(z)=\sum_j(w_jv_j^{(2)}+a_jv_j)z^j,\qquad
 C(z)=\sum_jw_jv_j^2z^j.
\]
We have $B_2(0)=0$, $\deg C\le N-1$, and
\begin{equation}\label{eq:lead}
 [z^N]B_2=2^{-N}A_N,
\end{equation}
since $v_N=0$ and $v_N^{(2)}=A_N$.

Expanding polynomial divisibility one order further gives
\begin{align*}
 \Phi_2(1+y)&\equiv-\frac{\Phi_1(1+y)}{y^3}-\frac{\Phi_0}{y^6},\\
 \Phi_2(-1+y)&\equiv\frac{\Phi_1(-1+y)}{y^3}-\frac{\Phi_0}{y^6}.
\end{align*}
Although individual displayed quotients can have negative powers for $N<6$, their indicated combinations are analytic: this follows directly from first-order divisibility. Subtraction gives
\[
 \ii B_2-yC\equiv-\frac{A+\ii yB}{y^3},
\]
or, in the $x$ coordinate,
\begin{equation}\label{eq:R}
 B_2\equiv R(x):=\frac{\pi^3A}{\ell^3}
                +\frac{\pi^2B}{\ell^2}-\frac{\ell C}{\pi}\pmod{x^N}.
\end{equation}
Define $\mathcal S_N(R)=\sum_{k=0}^{N-1}[x^k]R(x)$. The degree-below-$N$ remainder in \eqref{eq:R}, evaluated at $z=0$, is $\mathcal S_N(R)$ because $x=1$ there. The condition $B_2(0)=0$ therefore fixes
\[
 B_2=\sum_{k=0}^{N-1}[x^k]R\,x^k-\mathcal S_N(R)x^N,
 \qquad A_N=-2^N\mathcal S_N(R).
\]

For the $A,B$ terms in \eqref{eq:R}, substitution of \eqref{eq:A}--\eqref{eq:B} factors out $\pi^5 2^{-N}x^{N-5}$ and leaves
\[
 -\frac{(x-2)^2x^3}{4\ell^3}
 +\frac{(x-1)(2-x)x^2}{2\ell^2}
 =\frac{x^2}{6}-\frac{x^3}{4}+\frac{11x^4}{120}+O(x^5).
\]
It follows, for every $N\ge3$, that
\begin{equation}\label{eq:SAB}
 \mathcal S_N(R_{AB})=\frac{\pi^5}{120\,2^N}.
\end{equation}

For the final term write $c_j=w_jv_j^2$, so $c_0=c_N=0$ and $C(z)=\sum_{j=1}^{N-1}c_jz^j$. Since
$C(x-1)/(1-x)=-\sum_jc_j(x-1)^{j-1}$,
\begin{align}
 \mathcal S_N(-\ell C/\pi)
 &=\frac1\pi\sum_{j=1}^{N-1}c_j[x^{N-1}]\ell(x-1)^{j-1}\notag\\
 &=-\frac1\pi\sum_{j=1}^{N-1}c_j
           \frac{(N-j-1)!(j-1)!}{(N-1)!}\notag\\
 &=-\frac{N}{\pi 2^N}\sum_{j=1}^{N-1}\frac{v_j^2}{j(N-j)}.\label{eq:SC}
\end{align}
The middle equality follows by expanding the logarithm, or integrating
$\int_0^1 t^{N-j-1}(1-t)^{j-1}\dd t$. The exact power-sum identity
\begin{equation}\label{eq:sum}
 \sum_{j=1}^{N-1}j(N-j)(N-2j)^2
   =\frac{N(N-1)(N-2)(N+1)(N+2)}{30}
\end{equation}
now gives, using \eqref{eq:v},
\[
 \mathcal S_N(R_C)=
 -\frac{\pi^5(N+1)(N+2)}{120\,2^N(N-1)(N-2)}.
\]
Combining with \eqref{eq:SAB},
\begin{equation}\label{eq:AN}
 \boxed{A_N=\frac{N\pi^5}{20(N-1)(N-2)}.}
\end{equation}
This calculation uses finite power-series coefficients and polynomial identities only. An exact symbolic verification is supplied with the paper.

\section{Completion, examples and the prime-error application}
\begin{proof}[Proof of Theorem~\ref{thm:main}]
The local branch and uniqueness follow from Proposition~\ref{prop:branch}. Proposition~\ref{prop:ideal} and \eqref{eq:AN} give the normalized expansion. Since $e_1=0$, Newton's identity is $p_3=3e_3$. Substitution and \eqref{eq:scale} prove \eqref{eq:main}. For a centered direction $\delta_j=ta_j$ with $\sum a_j^3\ne0$, its excess is a strictly positive constant times $t^6+O(t^8)$, which proves optimality of the sixth-order exponent.
\end{proof}

\begin{proof}[Proof of Theorem~\ref{thm:equality}]
For $N=3,4$ there is just one odd coordinate after $e_1=0$, namely $e_3$. Equation~\eqref{eq:ideal} is therefore
$L-N\pi=e_3^2G(\boldsymbol e)$, with $G(0)=A_N>0$. Shrinking the neighborhood makes $G>0$. Thus equality is equivalent to $e_3=0$. The centered cubic is $y^3+e_2y-e_3$, and the centered quartic is $y^4+e_2y^2-e_3y+e_4$. With real roots, $e_3=0$ is exactly the asserted reflection symmetry. Rescale by $\kappa$.
\end{proof}

\begin{corollary}[An unequal-order two-frequency family]\label{cor:example}
For sufficiently small real $\epsilon$, prescribe order two at $1-\epsilon$ and order one at $1+2\epsilon$. Then
\[
 B(1-\epsilon,1-\epsilon,1+2\epsilon)
 =3\pi+\frac{3\pi^5}{10}\epsilon^6+O(\epsilon^8).
\]
For nonzero $\epsilon$ sufficiently close to zero the width is strictly larger than $3\pi$.
\end{corollary}
\begin{proof}
The mean is one and $p_3=6\epsilon^3$; apply the two theorems.
\end{proof}
For $\epsilon=0.01$, an independent high-precision solution of the three centered moment equations gives width excess approximately $9.1902306864\cdot10^{-11}$, while the leading term is $9.1805905436\cdot10^{-11}$. These floating-point values illustrate the asymptotic and are not part of its proof. The accompanying program also checks $N=3,\ldots,10$ with repeated and distinct targets.

\paragraph{Comparison with separate cancellation.}
Convolving one two-point filter per occurrence gives width
\[
 B_{\rm sep}=\pi\sum_j\frac1{\lambda_j}
 =\frac{N\pi}{\kappa}+\frac{\pi}{\kappa}\sum_j\delta_j^2+O_N(d^3/\kappa).
\]
The optimum has no quadratic penalty. Its first possible penalty is sixth order and depends on the asymmetric part of the cluster; reflection-symmetric clusters retain the exact mean-based width. Thus the result identifies a structural saving rather than a small adjustment of a numerical constant.

\paragraph{Application to the repository's prime averages.}
Suppose a normalized prime error contains a finite group of modes
$e^{\ii\lambda_j y}$, with the critical-line Fourier description supplied by the repository's Dirichlet GRH input. Averaging $E(y+u)$ against the minimizing positive measure multiplies each such mode by $\Phi_\mu(\lambda_j)$ and cancels it. Prescribing order $m$ also cancels
$y^k e^{\ii\lambda y}$ for $0\le k<m$, and gives the elementary detuning estimate
\[
 |\Phi_\mu(\lambda+\eta)|\le\frac{(B|\eta|)^m}{m!}
\]
when the measure is anchored in $[0,B]$. This is Taylor's integral remainder and $|e^{\ii\omega u}|=1$ for real $\omega$.
In the original multiplicative scale, the required largest-to-smallest scale ratio is $e^B$.

For unnormalized powers $x^{\beta+\ii\lambda}$ with common real part $\beta$, replace the Fourier measure by
\[
 \dd\nu(u)=\frac{e^{-\beta u}\dd\mu(u)}{\int e^{-\beta v}\dd\mu(v)}.
\]
Its Laplace transform vanishes at $\beta+\ii\lambda$ to the same orders, and its support diameter is unchanged. Tilting back is the inverse correspondence, so the minimum width is unchanged. The filter theorem itself does not use RH or GRH; only this arithmetic identification of modes uses that input.

\paragraph{Scope and remaining questions.}
The constants $\epsilon_N,C_N$ are existence constants from the analytic argument, not explicit certified radii. The theorem is a fixed-$N$ local statement, with no uniform claim as $N$ grows. It does not settle the complete global phase diagram or the sign of the first nonzero term in every higher-dimensional direction with $p_3=0$. Computing the quadratic form in $e_3,e_5,\ldots$ beyond its first entry would address the latter question. The accompanying calculation is reproducible symbolic and numerical evidence; no Lean formalization is claimed.

\section*{Reproducibility and source record}
\texttt{check\_spectral\_skewness.py} checks the reciprocal-log expansion, the binomial first-variation formulas, the beta-integral coefficient identity, the polynomial power sum, and the final coefficient with exact SymPy arithmetic. Its optional numerical mode uses Hermite divided differences and high-precision symmetric-atom equations. \path{spectral_skewness_verification.json} records the results. The source comparison is pinned to \texttt{Article/sections/fil.tex}, Git blob
\texttt{e5528541d9708699556049c956ffefe9f4621924}, at the repository commit cited above; the relevant source labels are \texttt{fil:thm:split}, \texttt{fil:thm:differential}, \texttt{fil:cor:cluster}, and \texttt{fil:sec:open}.

\begin{thebibliography}{9}
\bibitem{ProveIt} V. Reshetnikov, \emph{ProveIt: Beyond the Riemann hypothesis for Dirichlet $L$-functions}, positive-filter section, revision \texttt{00cadb961bbe77c74e5d082c4948f8884d58a970}, consulted 27 November 2026.
\href{https://github.com/VladimirReshetnikov/ProveIt/blob/00cadb961bbe77c74e5d082c4948f8884d58a970/Analysis/GeneralizedRiemannHypothesis/Article/sections/fil.tex}{Pinned source file}.
\bibitem{Crain1996} E. Crain, W. Singhose and W. Seering, Derivation and properties of convolved and simultaneous two-mode input shapers, \emph{IFAC Proceedings Volumes} \textbf{29} (1996), no.~1, 7278--7283.
\href{https://doi.org/10.1016/S1474-6670(17)58856-4}{doi:10.1016/S1474-6670(17)58856-4}.
\bibitem{Singhose1997} W. Singhose, E. Crain and W. Seering, Convolved and simultaneous two-mode input shapers, \emph{IEE Proceedings---Control Theory and Applications} \textbf{144} (1997), no.~6, 515--520.
\href{https://doi.org/10.1049/ip-cta:19971439}{doi:10.1049/ip-cta:19971439}.
\end{thebibliography}
\end{document}
